\section{Climbing Stairs}
\label{sec:dp-climbing-stairs}

\problemstatement{A staircase has $n$ steps. From any step, you may
climb either $1$ or $2$ steps at a time. Count the number of distinct
ways to reach the top (step $n$), starting from step $0$.}

\begin{patternbox}
\emph{``Number of distinct ways to reach step $n$, taking $1$ or $2$
steps at a time''} is the direct counting analogue of
Section~\ref{sec:dp-fibonacci}'s recurrence: the last move to reach step
$n$ was either a single step from $n-1$, or a double step from $n-2$ --
so every way of reaching $n$ is exactly one way of reaching $n-1$,
extended by one step, \emph{or} one way of reaching $n-2$, extended by
two steps, and these two groups are disjoint and exhaustive.
\end{patternbox}

\subsection*{Deriving the recurrence}

Let $\textit{ways}(i)$ be the number of distinct ways to reach step $i$.
Since the only moves are $+1$ and $+2$, the last move into step $i$ came
from step $i-1$ or step $i-2$ -- so $\textit{ways}(i) =
\textit{ways}(i-1) + \textit{ways}(i-2)$. Base cases:
$\textit{ways}(0)=1$ (one way to be at the start: do nothing), and
$\textit{ways}(1)=1$ (only one way to reach step $1$: a single step from
$0$). This is exactly Fibonacci's recurrence, shifted by one index, with
different base cases.

\subsection*{Approach 1: Recursion}

\begin{lstlisting}
#include <bits/stdc++.h>
using namespace std;

int climbStairsRecursive(int n) {
    if (n <= 1) return 1;
    return climbStairsRecursive(n - 1) + climbStairsRecursive(n - 2);
}

int main() {
    cout << climbStairsRecursive(5) << endl; // 8
    return 0;
}
\end{lstlisting}

\begin{complexitybox}[Approach 1 Complexity]
\textbf{Time.} $O(2^n)$ -- identical overlapping-subproblem blowup as
plain Fibonacci.
\textbf{Space.} $O(n)$ for the recursion stack.
\end{complexitybox}

\subsection*{Approach 2: Memoization}

\begin{lstlisting}
#include <bits/stdc++.h>
using namespace std;

int climbStairsMemo(int n, vector<int>& memo) {
    if (n <= 1) return 1;
    if (memo[n] != -1) return memo[n];

    memo[n] = climbStairsMemo(n - 1, memo) + climbStairsMemo(n - 2, memo);
    return memo[n];
}

int main() {
    int n = 5;
    vector<int> memo(n + 1, -1);
    cout << climbStairsMemo(n, memo) << endl; // 8
    return 0;
}
\end{lstlisting}

\begin{complexitybox}[Approach 2 Complexity]
\textbf{Time.} $O(n)$.
\textbf{Space.} $O(n)$ memo array $+$ $O(n)$ recursion stack.
\end{complexitybox}

\subsection*{Approach 3: Tabulation}

\begin{lstlisting}
#include <bits/stdc++.h>
using namespace std;

int climbStairsTabulation(int n) {
    if (n <= 1) return 1;

    vector<int> dp(n + 1);
    dp[0] = 1;
    dp[1] = 1;
    for (int i = 2; i <= n; i++) {
        dp[i] = dp[i - 1] + dp[i - 2];
    }
    return dp[n];
}

int main() {
    cout << climbStairsTabulation(5) << endl; // 8
    return 0;
}
\end{lstlisting}

\subsection*{Dry run of the tabulation table}

\dryrun{Take $n=5$.}

\begin{center}
\begin{tabular}{c|c|c|c|c|c|c}
$i$ & 0 & 1 & 2 & 3 & 4 & 5 \\ \hline
$\textit{dp}[i]$ & 1 & 1 & 2 & 3 & 5 & 8
\end{tabular}
\end{center}

$\textit{dp}[5]=8$: indeed, there are $8$ distinct ways to climb $5$
steps using moves of size $1$ or $2$ (e.g.\ $1{+}1{+}1{+}1{+}1$,
$2{+}1{+}1{+}1$ in every position, $2{+}2{+}1$ in every order, and so
on).

\begin{complexitybox}[Approach 3 Complexity]
\textbf{Time.} $O(n)$.
\textbf{Space.} $O(n)$ for the table, no recursion stack.
\end{complexitybox}

\subsection*{Approach 4: Space Optimization}

\begin{lstlisting}
#include <bits/stdc++.h>
using namespace std;

int climbStairsSpaceOptimized(int n) {
    if (n <= 1) return 1;

    int prev2 = 1, prev1 = 1;
    for (int i = 2; i <= n; i++) {
        int current = prev1 + prev2;
        prev2 = prev1;
        prev1 = current;
    }
    return prev1;
}

int main() {
    cout << climbStairsSpaceOptimized(5) << endl; // 8
    return 0;
}
\end{lstlisting}

\begin{complexitybox}[Approach 4 Complexity]
\textbf{Time.} $O(n)$.
\textbf{Space.} $O(1)$.
\end{complexitybox}

\begin{takeawaybox}
Recognizing that ``count the ways to reach $n$ via moves from a small
fixed set'' always produces a recurrence that \textbf{sums} the way-counts
of the positions those moves originate from is the key insight -- the
specific move set (here, $\{1,2\}$) only determines \emph{how many}
earlier terms the recurrence sums, and which ones. Section~\ref{sec:dp-frog-jump}
generalizes this from counting ways to minimizing cost, with the sum
replaced by a minimum.
\end{takeawaybox}
